\chapter{The Meridian Companion Repository}
\label{ch:companion}

Every listing in this book was extracted from here. Every number was produced
here. It has no third-party runtime dependencies, which means you can clone it
and run everything with a stock Python.

\begin{center}
\texttt{https://github.com/krimler/Building-Awesome-MCP-Apps}
\end{center}

\section{Two minutes}

\begin{sh}
git clone https://github.com/krimler/Building-Awesome-MCP-Apps
cd Building-Awesome-MCP-Apps

make test      # 196 tests, about two seconds
make bench     # every measurement printed in this book
make lint      # prose rules and cross-reference integrity
make book      # build the PDF
\end{sh}

Python 3.11 or later. Nothing else for the code. The book build additionally
wants a TeX distribution and, for the figures, matplotlib.

\section{Layout}

\begin{plain}
meridian/
  protocol/          3,892 lines. A complete 2026-07-28 implementation
    jsonrpc.py         framing, resultType, SSE parsing
    meta.py            the _meta envelope, capabilities, RequestContext
    errors.py          the error-code allocation
    server.py          routing, primitives, schema validation
    client.py          the MRTR loop, caching, era detection
    stdio.py           subprocess transport, both sides
    http.py            Streamable HTTP, header validation, SSE
    inproc.py          the control group for benchmarks
    mrtr.py            input_required, elicitation, StateSealer
    subscriptions.py   subscriptions/listen
    cache.py           ttlMs and cacheScope
    tasks.py           the Tasks extension
    legacy.py          the dual-era bridge
  servers/           the four Meridian servers, plus fixtures
  host/              Host, AgentLoop, StubModel
  bench/             the measurement harness and results.json
  tests/             196 tests across five files
  serve.py           the launcher
\end{plain}

\section{Running the servers}

\begin{sh}
python3 -m meridian.serve risk                 # dual-era stdio
python3 -m meridian.serve risk --modern-only   # refuses the handshake
python3 -m meridian.serve risk --http 8931     # Streamable HTTP
python3 -m meridian.serve all --http 8931      # all four, sequential ports
\end{sh}

Dual-era is the default, because that is what connects to the clients shipping
in 2026 (\chapref{architecture}).

\subsection{The scenario flags}

\begin{table}[htbp]
\centering
\small
\begin{tabularx}{\textwidth}{@{}lXl@{}}
\toprule
\thead{Flag} & \thead{What it does} & \thead{Chapter} \\
\midrule
\texttt{-{}-fat} & Serves the 38-tool REST mirror & \ref{ch:tools} \\
\texttt{-{}-poisoned} & Plants the injection in the guidance corpus & \ref{ch:threats} \\
\texttt{-{}-modern-only} & Refuses \texttt{initialize} & \ref{ch:architecture} \\
\bottomrule
\end{tabularx}
\caption{Each flag reproduces a specific measurement or attack from the book.}
\label{tab:meridian-flags}
\end{table}

\section{With Claude Code}

The repository ships a \texttt{.mcp.json}:

\begin{sh}
claude -p "Assess ACC-1042 for risk, screen it for fraud, and get the 1Y
and 5Y reference rates." --mcp-config .mcp.json
\end{sh}

\begin{plain}
RISK:  55.4 elevated
FRAUD: watch 1
CURVE: 1Y=3.96 5Y=3.68
\end{plain}

Those values come straight from the fixtures. The full transcript, and what is
\emph{not} verifiable this way and why, is in
\texttt{meridian/VERIFICATION.md}.

\section{The four servers}

\begin{table}[htbp]
\centering
\small
\begin{tabularx}{\textwidth}{@{}lXX@{}}
\toprule
\thead{Server} & \thead{Demonstrates} & \thead{Chapters} \\
\midrule
\texttt{risk}
  & Task-shaped tools, MRTR approvals, Tasks, an App template, TTL'd resources
  & \ref{ch:tools}, \ref{ch:mrtr}, \ref{ch:tasks}, \ref{ch:apps} \\
\texttt{compliance}
  & Human-in-the-loop, audit logging, the injection payload
  & \ref{ch:mrtr}, \ref{ch:threats}, \ref{ch:access-control} \\
\texttt{fraud}
  & Fast read-only calls, runtime \texttt{listChanged}
  & \ref{ch:tools}, \ref{ch:resources} \\
\texttt{marketdata}
  & \texttt{cacheScope: public}, pagination
  & \ref{ch:resources} \\
\bottomrule
\end{tabularx}
\caption{Each server exists to carry specific arguments in the book.}
\label{tab:meridian-servers}
\end{table}

\section{The benchmark harness}

\begin{sh}
make bench                                       # everything
python3 -m meridian.bench.run cache fanout       # named scenarios
python3 -m meridian.bench.run --list             # scenario names
python3 -m meridian.bench.run --json out.json    # machine-readable
\end{sh}

\begin{table}[htbp]
\centering
\small
\begin{tabularx}{\textwidth}{@{}lXl@{}}
\toprule
\thead{Scenario} & \thead{Answers} & \thead{Chapter} \\
\midrule
\texttt{transport}     & What does each transport cost, minus execution? & \ref{ch:transport} \\
\texttt{coldstart}     & Spawn per call, or warm pool? & \ref{ch:transport} \\
\texttt{catalogue}     & What does catalogue bloat cost per turn? & \ref{ch:tools} \\
\texttt{cache}         & What is honouring \texttt{ttlMs} worth? & \ref{ch:resources} \\
\texttt{fanout}        & What does parallelism buy, at what latency? & \ref{ch:tools} \\
\texttt{mrtr}          & What does one elicitation really cost? & \ref{ch:mrtr} \\
\texttt{loop}          & Where does a whole task's time go? & \ref{ch:measuring} \\
\texttt{optimization}  & The full pass, attributed change by change & \ref{ch:optimizing} \\
\texttt{serialisation} & How much of a request is envelope? & \ref{ch:architecture} \\
\bottomrule
\end{tabularx}
\caption{The canonical run lives in \texttt{meridian/bench/results.json}, and
every figure in this book that plots data reads it.}
\label{tab:bench-reference}
\end{table}

\subsection{What is simulated}

Stated here as well as at every point of use.

Model inference latency and token pricing are \textbf{modelled}, from the
constants in \texttt{meridian/host/model.py}:

\begin{py}
TTFT_MEAN_MS         = 340.0
TTFT_STDEV_MS        = 90.0
MS_PER_OUTPUT_TOKEN  = 7.5
CHARS_PER_TOKEN      = 4.0
INPUT_USD_PER_MTOK   = 3.00
OUTPUT_USD_PER_MTOK  = 15.00
\end{py}

Everything else is \textbf{measured}: serialisation, transport, server execution,
cache behaviour, round-trip counts, catalogue token cost, process cold starts.

Transport figures are loopback figures, deliberately, so protocol overhead is
visible without a network in the way. Chapter~\ref{ch:measuring} explains how to
add your own.

Edit those constants, rerun \texttt{make bench}, and every number in the book
moves with you.

\section{The test suite}

\begin{table}[htbp]
\centering
\small
\begin{tabularx}{\textwidth}{@{}lrX@{}}
\toprule
\thead{File} & \thead{Tests} & \thead{Covers} \\
\midrule
\texttt{test\_protocol.py} & 57
  & Envelope, \texttt{resultType}, error codes, caching, pagination, schemas,
    MRTR, \texttt{StateSealer}, \texttt{x-mcp-header} \\
\texttt{test\_transports.py} & 28
  & Header encoding and validation, removed endpoints, SSE progress,
    subscriptions, real stdio subprocesses \\
\texttt{test\_integration.py} & 33
  & Host routing, consent, MRTR end to end, Tasks, caching, loop budgets \\
\texttt{test\_legacy.py} & 16
  & Era selection, legacy translation, bridge transparency \\
\bottomrule
\end{tabularx}
\caption{196 tests, roughly two seconds, no model required.}
\label{tab:test-files}
\end{table}

\section{Reading it in an evening}

If you want to understand the protocol from the code, this order works:

\begin{enumerate}[leftmargin=1.6em, itemsep=3pt]
\item \texttt{protocol/meta.py} \, The envelope. This is the 2026 change.
\item \texttt{protocol/jsonrpc.py} \, Framing, and \texttt{resultType}.
\item \texttt{protocol/server.py} \, \texttt{handle()} is the whole server.
\item \texttt{protocol/client.py} \, \texttt{call()} is the MRTR loop.
\item \texttt{protocol/mrtr.py} \, \texttt{StateSealer}, and why each binding
exists.
\item \texttt{protocol/http.py} \, Header validation. The security-critical part.
\item \texttt{servers/risk.py} \, All of it, applied.
\item \texttt{tests/test\_protocol.py} \, The specification, as assertions.
\end{enumerate}

Roughly two thousand lines for the first six, and they contain every mechanism
this book describes.

\section{Using it as a starting point}

You may. It is a teaching implementation, so before production:

\begin{itemize}[leftmargin=1.6em, itemsep=3pt]
\item Replace \texttt{TaskStore} with something shared and durable
(\chapref{tasks}).
\item Load the \texttt{StateSealer} secret from a secret manager, fleet-wide
(\chapref{mrtr}).
\item Add real authorization; the authenticator hook is a callable and does
nothing by default (\chapref{authorization}).
\item Replace \texttt{ThreadingHTTPServer} with a production server.
\item Add real observability. The built-in counters are deliberately crude
(\chapref{measuring}).
\item Consider an SDK instead. By the time you read this they will have caught
up with \texttt{2026-07-28}, and this exists to show you the wire rather than to
compete with them.
\end{itemize}

\section{Reproducing the book itself}

\begin{sh}
make figures    # regenerate every figure to PDF and SVG
make book       # the PDF
make site       # the website
make lint       # prose rules and reference integrity
make verify     # tests, benchmarks, and the Claude Code runs
\end{sh}

The figure pipeline builds each figure to \emph{both} PDF and SVG from one
source, so the printed book and the website cannot drift. The website is
generated from the same LaTeX as the PDF, so there is exactly one copy of every
sentence.

\section{Contributing}

Corrections welcome, and reproduction failures most of all. The entire premise
is that the numbers are checkable, so a number you cannot reproduce is the most
useful bug report available. Include your platform, your Python version, and the
output of \texttt{make bench}.
